\documentclass[10pt,a4paper]{amsart}
\usepackage[margin=19mm]{geometry}
\usepackage{amsmath,amssymb,mathtools}
\usepackage[scr=rsfs,cal=euler]{mathalfa}
\usepackage{xfrac}
\usepackage[unicode,hidelinks,bookmarks=false]{hyperref}
\newcommand{\ie}{i.e.\ }
\newcommand{\cf}{cf.\ }
\newcommand{\resp}{respectively\ }
\newcommand{\Subsection}{\subsection}
\providecommand{\nobreakdash}{\nobreak\hbox{-}}
\setlength{\emergencystretch}{2em}
\multlinegap0pt
\usepackage{bm}
\usepackage{bbm}
\usepackage{dsfont}
\usepackage{xspace}
\usepackage{cancel}
\usepackage{mathtools}
\usepackage{amsthm}
\makeatletter
\@ifundefined{theorem}{\newtheorem{theorem}{Theorem}}{}
\@ifundefined{lemma}{\newtheorem{lemma}[theorem]{Lemma}}{}
\makeatother
\newcommand{\dd}{\, \mathrm{d}}
\newcommand{\lzd}{{L^2(\Omega)^d}}
\newcommand{\lzo}{L^2(\Omega)}
\title[Determination of active forces in actomyosin systems as inve]{Determination of active forces in actomyosin systems as inverse source problems for the Stokes equation}
\author{Emily Klass, Tram Thi Ngoc Nguyen, Nilay Cicek, Yoav G. Pollack, Sarah K\"oster, Andreas Janshoff, Anne Wald}
\date{Source: arXiv:2601.09356v1 (CC BY 4.0)}
\begin{document}
\maketitle

\noindent\emph{Source and licence.} Determination of active forces in actomyosin systems as inverse source problems for the Stokes equation. Version arXiv:2601.09356v1 (CC BY 4.0), distributed under the Creative Commons Attribution 4.0 International licence, as recorded in the arXiv licence metadata for this exact version and shown on the abstract page https://arxiv.org/abs/2601.09356v1. Full rights notice: licenses/arxiv-2601.09356-CC-BY-4.0.txt.\par\smallskip

\noindent\textit{Editorial scope.} Counted unit: the Lemma whose source label is \texttt{lemma:banach\_robin} in the article (source0.tex lines 662--676, source label \texttt{lemma:banach\_robin}), delivered together with the article's own complete original proof of it (source0.tex lines 677--718, one \texttt{proof} environment of 42 lines). No mathematics is written, repaired or re-derived: statement and proof are the article's own text line for line. Required context retained uncounted: none. Every definition, display and label of those lines is retained unchanged. Omitted from this package and not redistributed: 1--661, 719--1169. Nothing inside a retained excerpt is omitted or altered: every retained body is delivered exactly as the article's lines are written, so no omission receipt of any kind accompanies this package. Citations: the retained excerpts carry no citation command at all, so no bibliography entry of the article is transcribed and no reference is redistributed. Reference adaptation: none was needed and none was made; every reference inside the delivered unit resolves inside it, so no unresolved reference or citation is produced. Printed slips kept literal: no printed word, symbol, hypothesis or number of the counted statement or of its proof was altered. Layout and class: the article's own class is \texttt{scrarticle} with packages including amsmath, amssymb, amsthm, float, comment, graphicx, textcomp, appendix, wrapfig, subcaption, srcltx, hyperref, algorithm, algpseudocode; the wrapper is the lane's pinned \texttt{amsart} editorial wrapper at 10pt on A4, which restates the theorem-like environments the retained text needs and adds the article's own macro definitions that the retained text needs (\texttt{\textbackslash dd}, \texttt{\textbackslash lzd}, \texttt{\textbackslash lzo}), with extra wrapper packages (bm, bbm, dsfont, xspace, cancel, mathtools). The delivered numbering is the unit's own. Assets: no retained span contains a figure environment, a graphics-inclusion command, a PDF-page inclusion, a table or a TikZ picture; no image file is copied into this package, the package declares no asset dependency and no image file is redistributed with it. Licence: version arXiv:2601.09356v1 of this article is distributed under the Creative Commons Attribution 4.0 International licence, as recorded in the arXiv licence metadata for this exact version and shown on the abstract page https://arxiv.org/abs/2601.09356v1. A full rights notice is delivered at licenses/arxiv-2601.09356-CC-BY-4.0.txt. Characters: every retained excerpt is pure ASCII, so the delivered engine, XeLaTeX with its default Latin Modern text font, reports no missing character.
\begin{lemma}[Banach space adjoint of $S_R$]
\label{lemma:banach_robin}
    The Banach adjoint operator of $S_R$ is given by
    \begin{displaymath}
        S_R^{\#}:\left(\lzd \times \lzo\right) \rightarrow \lzd, \quad (\mathbf{w},w_p) \mapsto \boldsymbol{\varphi},
    \end{displaymath} where $\left(\boldsymbol{\varphi},\varphi_p\right)\in \left(\lzd \times \lzo\right)$ solves the adjoint equation
      \begin{equation} \label{eq-adjoint}
 \begin{split}
  -\nabla \cdot \big( \eta D(\boldsymbol{\varphi}) \big) - \nabla \varphi_p &= \mathbf{w} \quad \text{in } \Omega\\
  \nabla \cdot \boldsymbol{\varphi} &= w_p \quad \text{in } \Omega\\
  \boldsymbol{\varphi}_{\bot} &= \mathbf{0} \quad \text{on } \partial\Omega \\
  \big(\eta (\nabla \boldsymbol{\varphi})\mathbf{n} \big)_{\parallel} &= -\lambda \boldsymbol{\varphi}_{\parallel}\quad \text{on } \partial\Omega.
  \end{split}
\end{equation}
\end{lemma}

\begin{proof}  Take any $(\mathbf{w},w_p)\in \lzd \times \lzo$ and any $\mathbf{h}\in\lzd$. 
    Let $(\mathbf{v},p):=S_R(\mathbf{h})$ and $\left(\boldsymbol{\varphi},\varphi_p\right)$ as in \eqref{eq-adjoint}.
    Our aim is to show that the following identity holds
    \begin{equation}
    \langle (\mathbf{w},w_p),S_R(\mathbf{h})\rangle = 
        \langle \mathbf{w}, \mathbf{v} \rangle_{\lzd} + \langle w_p , p \rangle_{\lzo} = \langle \boldsymbol{\varphi} , \mathbf{h} \rangle =
        \langle S_R^{\#} (\mathbf{w},w_p), \mathbf{h} \rangle_{\lzd}.
        \label{eq:adj_cond}
    \end{equation}
    We proceed by rewriting the right-hand side of \eqref{eq:adj_cond}. Firstly, integration-by-part yields
    \begin{equation}
    \begin{split}
         \langle \boldsymbol{\varphi}, - \nabla \cdot (\eta D(\mathbf{v})) \rangle &= \int_\Omega \eta \nabla \boldsymbol{\varphi} : D(\mathbf{v}) \dd \mathbf{x} - \int_{\partial \Omega} \boldsymbol{\varphi}^{T} (\eta D(\mathbf{v}) \mathbf{n} ) \dd \boldsymbol{\sigma} \\
         & = \int_\Omega \eta D(\boldsymbol{\varphi}) : D(\mathbf{v}) \dd \mathbf{x} + \int_{\partial\Omega} \lambda \boldsymbol{\varphi}_{||}^T \mathbf{v}_{||} \dd \boldsymbol{\sigma} \\
         & = \langle \mathbf{v}, - \nabla \cdot (\eta D(\boldsymbol{\varphi}) )\rangle,
         \label{eq:srad_proof_v}
    \end{split}
    \end{equation}
    where the second equality is obtained with symmetries of $D(\mathbf{v})$ and $D(\boldsymbol{\varphi})$.
    Secondly, using the divergence theorem we have 
    \begin{equation}\label{eq:srad_proof_p}
        \begin{split}
            \langle \boldsymbol{\varphi}, \nabla p\rangle = - \int_\Omega  p (\nabla \cdot\boldsymbol{\varphi}) \dd \mathbf{x} + \int_{\partial\Omega} \boldsymbol{\varphi}^T (p \mathbf{n}) \dd \boldsymbol{\sigma} 
                                = \langle - \nabla \cdot \boldsymbol{\varphi}, p\rangle,
        \end{split}
    \end{equation}
    where $\int_{\partial\Omega} \boldsymbol{\varphi}^T (p \mathbf{n}) \dd \boldsymbol{\sigma} = 0$ as $\boldsymbol{\varphi}_\bot = \mathbf{0}$.
    The incompressibility of the velocity $\mathbf{v}$ allows us to introduce $\varphi_p\in\lzo$ and together with the divergence theorem gives
    \begin{equation}\label{eq:srad_proof_incomp}
     0=(\nabla \cdot \mathbf{v} , \varphi_p)= -(\mathbf{v}, \nabla \varphi_p) .
    \end{equation}
    Summing up \eqref{eq:srad_proof_v}, \eqref{eq:srad_proof_p} and \eqref{eq:srad_proof_incomp} leads to
    \begin{equation}
        \begin{split}
            \langle \boldsymbol{\varphi},\mathbf{h}\rangle 
            &= \langle \boldsymbol{\varphi}, - \nabla \cdot (\eta D(\mathbf{v})) - \nabla p \rangle 
            = \langle \mathbf{v}, - \nabla \cdot (\eta D(\boldsymbol{\varphi}) ) - \nabla \boldsymbol{\varphi} \rangle + \langle \nabla \cdot \boldsymbol{\varphi} , p \rangle \\
            & = \langle \mathbf{v},\mathbf{w}\rangle + \langle p, w_p \rangle.
        \end{split}
    \end{equation}
which is the left-hand side of \eqref{eq:adj_cond}, and the proof is complete.
\end{proof}

\end{document}
